%==============================================================================
%  STAT 5200 - Applied Econometrics I
%  Homework 6
%
%  Compile with:  pdflatex "HW 6 - Report.tex"  (run twice for refs)
%
%  Directory layout assumed by this file:
%     HW 6 - Report.tex
%     R/        (R scripts)
%     Figures/  (generated figures)
%     Output/   (captured console output)
%==============================================================================
\documentclass[11pt]{article}

\usepackage[margin=1in]{geometry}
\usepackage[T1]{fontenc}
\usepackage[utf8]{inputenc}
\usepackage{amsmath,amssymb}
\usepackage{booktabs}
\usepackage{graphicx}
\usepackage{xcolor}
\usepackage{listings}
\usepackage[font=small,labelfont=bf]{caption}
\usepackage[hidelinks]{hyperref}

%------------------------------------------------------------------------------
%  Code presentation
%
%  Code is set as selectable text on a fixed character grid, source and console
%  output are given visually distinct styles, and source lines are kept to 76
%  characters so that nothing wraps.
%------------------------------------------------------------------------------
\definecolor{codekeyword}{RGB}{0,86,148}
\definecolor{codestring} {RGB}{163,21,21}
\definecolor{codecomment}{RGB}{106,115,125}
\definecolor{codebg}     {RGB}{248,248,250}
\definecolor{coderule}   {RGB}{220,220,226}

\lstdefinestyle{Rcode}{
  language        = R,
  basicstyle      = \ttfamily\footnotesize,
  keywordstyle    = \color{codekeyword}\bfseries,
  stringstyle     = \color{codestring},
  commentstyle    = \color{codecomment}\itshape,
  numbers         = left,
  numberstyle     = \ttfamily\scriptsize\color{codecomment},
  numbersep       = 9pt,
  backgroundcolor = \color{codebg},
  frame           = single,
  rulecolor       = \color{coderule},
  framesep        = 4pt,
  xleftmargin     = 15pt,
  xrightmargin    = 0pt,
  columns         = fixed,
  basewidth       = 0.5em,
  keepspaces      = true,
  tabsize         = 2,
  showstringspaces = false,
  breaklines      = true,
  breakatwhitespace = true,
  postbreak       = \mbox{\textcolor{codekeyword}{$\hookrightarrow$}\space},
  captionpos      = b,
  abovecaptionskip = 6pt,
}

\lstdefinestyle{Routput}{
  language        = {},
  basicstyle      = \ttfamily\footnotesize,
  keywordstyle    = \ttfamily,
  commentstyle    = \ttfamily,
  stringstyle     = \ttfamily,
  numbers         = none,
  backgroundcolor = \color{white},
  frame           = leftline,
  rulecolor       = \color{coderule},
  framesep        = 6pt,
  xleftmargin     = 20pt,
  breaklines      = true,
  captionpos      = b,
  abovecaptionskip = 6pt,
}

\lstset{style=Rcode}
\renewcommand{\lstlistingname}{Listing}

% Short inline code fragments.
\newcommand{\code}[1]{\texttt{#1}}

% Operators used throughout.
\DeclareMathOperator{\Var}{Var}
\DeclareMathOperator{\Cov}{Cov}
\DeclareMathOperator{\Avar}{Avar}
\DeclareMathOperator{\sd}{sd}
\DeclareMathOperator*{\plim}{plim}

% Placeholder marker. Every occurrence should be gone before submission;
% search the source for "\todo" to find what is still outstanding.
\newcommand{\todo}[1]{\textcolor{red}{\textbf{[TODO:} #1\textbf{]}}}

%------------------------------------------------------------------------------
\title{\vspace{-2em}Homework 6\\[0.3em]
       \large STAT 5200 --- Applied Econometrics I --- Fall 2026}
\author{Apostolos Kyratzis - 13130833}
\date{\today}

\begin{document}
\maketitle

\clearpage
\section{Q1}
\label{sec:q1}
\subsection{(1)}
\todo{Compare your HW5 answers with the solutions and list any missed questions.}

\subsection{(2)}
\todo{Discuss what you could have done better.}

\clearpage
\section{Q2}
\label{sec:q2}
Starting from the population wage equation
\[
    y=\beta_0+\beta_1s+u,
\]
take the covariance of both sides with the instrument $z$. By linearity of
covariance,
\begin{align*}
    \Cov(y,z)
    &=\Cov(\beta_0,z)+\beta_1\Cov(s,z)+\Cov(u,z)\\
    &=\beta_1\Cov(s,z).
\end{align*}
The covariance with the constant $\beta_0$ is zero, and instrument validity
implies $\Cov(u,z)=0$. Relevance implies $\Cov(s,z)\ne0$, so division gives
\[
    \boxed{\displaystyle \beta_1=\frac{\Cov(y,z)}{\Cov(s,z)}.}
\]
Thus, although schooling is correlated with the wage equation's error, a valid
instrument identifies its slope through variation in schooling that is
uncorrelated with that error. This argument assumes the relevant second moments
exist and follows Notes 8.

\clearpage
\section{Q3}
\label{sec:q3}
We use the regression results reported in the Lab07 notes.

\subsection{(1)}
The estimated coefficient on education is $0.10783$ under OLS, $0.04653$
using mother's education as the sole excluded instrument, and $0.07913$
using mother's, father's, and husband's education as three excluded instruments.
The following table reports the coefficients and their standard errors.
\begin{center}
\begin{tabular}{lrrr}
\toprule
Estimator & Estimate & Conventional SE & HC1 robust SE\\
\midrule
OLS & 0.10783 & 0.01440 & 0.01362\\
IV: mother's education & 0.04653 & 0.03800 & 0.03880\\
2SLS: three instruments & 0.07913 & 0.02241 & 0.02228\\
\bottomrule
\end{tabular}
\end{center}
The OLS coefficient on education is $0.1078$, compared with $0.0791$
using three-instrument 2SLS. Assuming the instruments are valid, this
difference suggests that education is endogenous and that OLS overestimates
its effect on log wages.

\subsection{(2)}
The lower 2SLS estimate suggests that OLS overestimates the effect of
education on log wages. This direction is plausible if, after accounting
for the included controls, people with more schooling also have unobserved
characteristics, such as ability or motivation, that independently increase
their wages. OLS then attributes part of those characteristics' effects to
education. Assuming the instruments are valid, 2SLS isolates variation in
schooling unrelated to these unobserved wage determinants, giving a lower
estimate. Ability and motivation are possible explanations for the upward
bias; the estimates do not establish which mechanism is responsible.

\subsection{(3)}
Lab 7 reports an OLS coefficient of $0.1078320$ on education and a
standard error of $0.0144021$. The 95\% confidence interval is
\[
    \hat\beta_{\mathrm{educ}}\pm1.96\,\mathrm{SE}(\hat\beta_{\mathrm{educ}})
    =0.1078320\pm1.96(0.0144021)
    =\boxed{[0.07960,\ 0.13606]}.
\]

\subsection{(4)}
Using mother's, father's, and husband's education as instruments, Lab 7
reports a 2SLS coefficient of $0.0791330188$ on education and a standard
error of $0.0224050892$. The 95\% confidence interval is
\[
    \hat\beta_{\mathrm{educ}}\pm1.96\,\mathrm{SE}(\hat\beta_{\mathrm{educ}})
    =0.0791330188\pm1.96(0.0224050892)
    =\boxed{[0.03522,\ 0.12305]}.
\]

\subsection{(5)}
The conventional OLS interval has width $0.05646$, compared with $0.08783$
for 2SLS. Thus the 2SLS interval is about 56\% wider. OLS has a smaller
standard error because it uses more of the observed variation in education
to estimate the coefficient. 2SLS relies on the variation predicted by the
instruments, after accounting for the controls. Since the instruments
explain only part of the variation in education, the 2SLS estimate is less
precise. Nevertheless, a
narrower interval is not always preferable. If education is endogenous, OLS
can estimate the wrong population quantity very precisely, so its narrower
interval can be misleading about the structural effect. Under valid and
relevant instruments, 2SLS addresses this endogeneity at a cost in precision.
Instrument validity and appropriate coverage matter alongside interval width.

\clearpage
\section{Q4}
\label{sec:q4}
\subsection{(1)}
I loaded the Mroz dataset, kept the 428 working women, and executed the
provided code.
\begin{lstlisting}[style=Rcode]
# Run from the Homework 6 folder.
mroz <- read.csv("Data/mroz.csv")
mroz <- subset(mroz, inlf == 1)
mroz$logwage <- mroz$lwage
mroz$exper2 <- mroz$expersq

with(mroz, {
# --- 1st stage
mm_1 = lm(educ ~ exper + exper2 + age + kidslt6 + kidsge6 + motheduc + fatheduc + huseduc);
educ_hat = matrix(mm_1$fitted.values); # xhat
# --- 2nd stage
mm2 = lm(logwage ~ exper + exper2 + age + kidslt6 + kidsge6 + educ_hat);
})
\end{lstlisting}

\subsection{(2)}
Yes, the point estimates match the three-instrument 2SLS estimates in
Lab 7. The coefficient on education is $0.07913301$.

Lab 7 uses matrix calculations. Following Notes 8, equations (18)--(20)
on p.\ 5, these are
\[
 \widehat X=Z(Z'Z)^{-1}Z'X,\qquad
 \widehat\beta_{\mathrm{2SLS}}
   =(\widehat X'X)^{-1}\widehat X'y
   =(\widehat X'\widehat X)^{-1}\widehat X'y.
\]
Here $y$ contains log wages, $X$ contains education and the controls,
and $Z$ contains the instruments and the controls; both matrices include
an intercept. The first calculation obtains fitted education and leaves
the controls unchanged. The second obtains the wage-regression coefficients.
Q4 performs the same steps using two \code{lm()} calls instead of matrix
calculations. Since the sample, instruments, and controls are the same,
the coefficients are the same. Notes 8, p.\ 10, describes this two-regression
implementation.

\subsection{(3)}
No, the standard errors differ. For education, Q4's second-stage
\code{lm()} gives $0.02342310$, compared with Lab 7's conventional 2SLS
standard error of $0.02240509$.

As explained in Notes 8, pp.\ 9--10, the correct 2SLS standard errors
use residuals calculated with the original regressors:
\[
 \widehat u=y-X\widehat\beta.
\]
Lab 7 calculates these residuals separately. In contrast, the second-stage
\code{lm()} calculates residuals using fitted education:
\[
 e=y-\widehat X\widehat\beta.
\]
The different residuals give different estimates of the error variance,
and hence different standard errors. The default standard errors from
the second-stage \code{lm()} are therefore not the correct 2SLS standard
errors; the correction is given in Notes 8, equations (25)--(26).

\clearpage
\section{Q5}
\label{sec:q5}
\subsection{(1)}
I selected Liu et al.\ (2026), published in
\emph{Manufacturing \& Service Operations Management}.
Section 6.5 uses instrumental variables to address endogeneity in the
relationship between workload and nurse attrition.

\subsection{(2)}
Liu, Blair (Lianlian), Diwas Singh KC, Bradley R. Staats, and Michael P.
Fundora. 2026. ``Frontiers in Operations: Operational Overload: The Impact
of Workload on High-Skilled Workforce Attrition.''
\emph{Manufacturing \& Service Operations Management} 28(5): 1458--1473.
\url{https://doi.org/10.1287/msom.2025.0037}.

\subsection{(3)}
The paper examines how different dimensions of workload affect ICU nurses'
voluntary decisions to leave their jobs, with the goal of improving
staffing and retention. It distinguishes responsibility within a shift,
overtime, exposure to patient deaths, and cumulative workload. The results
associate greater responsibility with lower attrition, while overtime and
exposure to patient deaths increase attrition. Cumulative workload has a
U-shaped relationship with attrition: moderate increases initially reduce
quitting, but excessive workload increases it. The authors also find that
coworker support can reduce some of the negative effects of workload.

\subsection{(4)}
The main model is a fixed-effects logit (Table 2, column 5):
\[
 \Pr(Y_{jt}=1\mid x_{jt})
 =\frac{\exp(x_{jt}'\beta)}{1+\exp(x_{jt}'\beta)}.
\]
Here $Y_{jt}=1$ if nurse $j$ voluntarily leaves four to eight weeks after
shift $t$. The regressors are nurse responsibility, overtime shifts,
emotional fatigue, cumulative workload and its square, plus four controls:
patient--nurse ratio, bed occupancy, unit experience, and current-shift
workload. This gives nine regressor terms, alongside nurse, month, and
year fixed effects. The estimation sample contains 17,606 nurse-shift
observations from 71 nurses who voluntarily left.

\subsection{(5)}
I focus on overtime shifts, which may be correlated with unobserved
reasons for quitting. The authors also consider nurse responsibility
and cumulative workload potentially endogenous.

\subsection{(6)}
A nurse already planning to leave may work extra shifts to save money
before resigning. The intention to quit then affects overtime, so the
observed relationship need not reflect only overtime causing attrition.

\subsection{(7)}
For overtime, the instrument is other ICU nurses' scheduled absences
during the previous week, excluding the focal nurse. These absences can
create staffing gaps and increase overtime opportunities; validity requires
that they affect quitting only through overtime, conditional on the controls.
The authors also propose lagged patient arrivals for nurse responsibility
and prior-month unscheduled peer absences for cumulative workload.

\subsection{(8)}
Section 6.5 presents first-stage equations and explains why the instruments
should predict workload. However, the main paper and online appendix do
not report estimated first-stage instrument coefficients or
excluded-instrument $F$ statistics, so numerical evidence of instrument
strength is not provided.

\subsection{(9)}
The baseline is logit rather than OLS; the IV correction uses residual
inclusion in the logit model. The coefficient on \emph{overtime shifts}
rises from $0.461$ (Table 2, column 5) to $54.28$ (Table 3). Assuming
valid instruments and a correctly specified model, this suggests downward
bias in the uncorrected overtime coefficient: ignoring endogeneity
understates overtime's positive effect on the log-odds of quitting.

In a logit model, the coefficient measures the change in
$\log[p/(1-p)]$ from one additional overtime shift, holding the other
regressors fixed, where $p$ is the probability of quitting and $p/(1-p)$
is the odds of quitting. Thus, $0.461$ does not mean a 46.1-percentage-point
increase in quitting probability. The probability change depends on the
starting probability and the other regressors, so the coefficient comparison
alone does not quantify how much the probability effect is understated.

\subsection{(10)}
Scheduled peer absences may leave a nurse with less help, less emotional
support, and greater stress during regular shifts, even without overtime.
These conditions could increase the likelihood of quitting independently
of overtime. If this channel remains after accounting for the controls,
the exclusion restriction fails and the instrument is invalid.

\clearpage
\appendix
\section{R code for Q4}
\label{app:iv-code}
Run the following code from the Homework 6 folder.
\begin{lstlisting}[style=Rcode,caption={The HW6 Q4 commands with data loading.},label={lst:q4-code}]
# Run from the Homework 6 folder.
mroz <- read.csv("Data/mroz.csv")
mroz <- subset(mroz, inlf == 1)
mroz$logwage <- mroz$lwage
mroz$exper2 <- mroz$expersq

with(mroz, {
# --- 1st stage
mm_1 = lm(educ ~ exper + exper2 + age + kidslt6 + kidsge6 + motheduc + fatheduc + huseduc);
educ_hat = matrix(mm_1$fitted.values); # xhat
# --- 2nd stage
mm2 = lm(logwage ~ exper + exper2 + age + kidslt6 + kidsge6 + educ_hat);
})
\end{lstlisting}

\section*{AI usage disclosure}
OpenAI Codex was used to help generate and adapt the R code included
in this report, using the course lab notes as references. Q3 uses the regression results printed in the Lab07 notes; the code
for Q4 was run to obtain its numerical comparisons.

\end{document}
